% =============================================================================
% Preambulo da variante TMLR (submissao anonima, double blind), xelatex.
%
% Entra via `pandoc -H`, portanto DEPOIS dos pacotes do template do pandoc
% (que ja carrega fontspec, amsmath, amssymb, longtable, booktabs, array e
% hyperref). Aqui fica so o que e do venue e o que e nosso.
%
% Por que `tmlr` e carregado aqui e nao num wrapper .tex com marcadores: o
% template default do pandoc ja emite exatamente a estrutura que o
% `tmlr.sty` espera — \title, \author, \maketitle, \begin{abstract} — e
% duplicar essa estrutura a mao significa mante-la em dois lugares.
%
% ⚠️ `\usepackage{tmlr}` SEM opcao = submissao anonima. Verificado no PDF, nao
% so no .sty: o nome real declarado no \author{} nao aparece no texto extraido,
% e sai "Anonymous authors / Paper under double-blind review". O build trata
% isso como pos-condicao e falha se o nome reaparecer.
%
% `tmlr.sty` faz \RequirePackage{natbib}. Este manuscrito carrega as
% referencias como footnotes escritas a mao, nao chaves BibTeX
% (`grep -c '[@' paper-tecnico-nox-mem.md` = 0), logo natbib e tmlr.bst ficam
% carregados e nao exercitados. Um camera-ready que converta as footnotes em
% BibTeX tem de usar `paper/tmlr/tmlr.bst`.
% =============================================================================

\usepackage{tmlr}

% ---------------------------------------------------------------------------
% Tabelas largas: o manuscrito tem tabelas de comparacao de 9 colunas, e a
% coluna de prosa colapsa para uma palavra por linha na largura do TMLR.
% Mesmo tratamento do preambulo do arXiv (paper/preamble.tex).
% ---------------------------------------------------------------------------
\usepackage{etoolbox}
\AtBeginEnvironment{longtable}{\footnotesize}
\setlength{\tabcolsep}{4pt}

% ---------------------------------------------------------------------------
% Margem: medido em 2026-09-11, a variante TMLR saia com 5 overfull \hbox, dois
% deles grandes (62,2pt e 54,5pt) — texto a invadir a margem. As duas causas
% sao (a) caminho absoluto longo em \texttt{}, que nao tem ponto de quebra, e
% (b) celula de tabela larga. `\emergencystretch` autoriza o TeX a esticar o
% espaco de uma linha antes de estourar a margem, e nao precisa de pacote.
%
% O remedio proprio para \texttt{} seria `hyphenat[htt]`, que NAO existe nesta
% TeX Live; depender dele tornaria o build nao-hermetico. O que sobrar de
% overfull se resolve no texto — um caminho absoluto de producao no meio da
% prosa ja e um cheiro antes de ser um problema de composicao.
%
% ⚠️ Conferir o numero de overfull depois de mexer aqui, e conferir ANTES que o
% build produziu paginas: na primeira tentativa o `hyphenat` ausente parou o
% LaTeX e o contador de overfull deu 0 — zero por nao ter medido, que se le
% igual a zero por estar limpo.
% ---------------------------------------------------------------------------
\emergencystretch=3em

% ---------------------------------------------------------------------------
% A CAUSA de todo glifo perdido: `fontspec` carregado e nunca instruido.
%
% O template do pandoc faz \usepackage{fontspec} sob XeTeX mas NAO chama
% \setmainfont, e sem isso o fontspec mantem a Latin Modern *Type1/EC*, de 256
% slots. Medido em 2026-09-11 neste manuscrito: 63 "Missing character" em
% ec-lmr10, ec-lmr8, ec-lmbx8, ec-lmtt8, ec-lmtt10 — os glifos alpha, seta,
% en dash, em dash, reticencias, menos, conjunto vazio, intersecao, aproximado
% e maior-ou-igual. Cada um **desaparece do PDF sem erro**: o build sai 0 e o
% caractere nao esta la.
%
% A Latin Modern OpenType traz todos e vem na TeX Live. Carregada pelo NOME DO
% ARQUIVO porque por nome de familia o fontspec nao a encontra nesta
% instalacao ("The font ... cannot be found").
%
% Isto substitui quatro remendos que eu tinha empilhado um a um, cada um a
% revelar o glifo seguinte: era UM defeito, nao cinco.
% ---------------------------------------------------------------------------
\setmainfont{lmroman10-regular.otf}[
  BoldFont       = lmroman10-bold.otf,
  ItalicFont     = lmroman10-italic.otf,
  BoldItalicFont = lmroman10-bolditalic.otf]
\setsansfont{lmsans10-regular.otf}[
  BoldFont   = lmsans10-bold.otf,
  ItalicFont = lmsans10-oblique.otf]
\setmonofont{lmmono10-regular.otf}[
  BoldFont   = lmmonolt10-bold.otf,
  ItalicFont = lmmono10-italic.otf,
  Scale      = MatchLowercase]

% ---------------------------------------------------------------------------
% Um glifo que nem a Latin Modern OpenType tem em modo texto: o alpha de
% "Mem-alpha", NOME de um sistema citado — soletra-lo seria errar o nome de
% quem se cita. Vem da Latin Modern Math, com `\char` porque o caractere e'
% ativo e reentraria.
%
% Os outros cinco (conjunto vazio, intersecao, uniao n-aria, aproximadamente,
% maior-ou-igual) foram reescritos em ASCII no fonte, e nao por preferencia: o
% construto `{\glifomat\char"..}` QUEBRA dentro de `\texttt{}` no documento
% completo ("Missing $ inserted"), embora compile isolado. Tres deles estavam
% exatamente la, num predicado de teoria de conjuntos num code span.
% ---------------------------------------------------------------------------
\newfontfamily\glifomat{latinmodern-math.otf}
\catcode`\α=\active \defα{{\glifomat\char"03B1}}
